\begin{textsample}{POS Dim 1 – rj – Score 25.00 – rj/rj\_ef\_ciencias\_humanas\_intro.txt}  \label{ex:f1_pos_254}
3.4.
A Área de Ciências Humanas Para o processo de \textbf{ensino-aprendizagem} em qualquer área do conhecimento, é cada vez mais necessário \textbf{que} não tenhamos \textbf{uma} compreensão do conhecimento como algo estanque e compartimentado, em \textbf{que} cada componente curricular ou área específica se basta em si mesma, fechando -se hermeticamente em \textbf{um} conjunto de conteúdo, \textbf{que}, sozinhos, já não podem dar conta de explicar \textbf{uma} realidade cada dia mais complexa e dinâmica, \textbf{que} \textbf{exige} das pessoas \textbf{um} pensar crítico e reflexivo acerca dos problemas concretos \textbf{que} a vida impõe.
Nesse sentido, qualquer proposta de ensino tem \textbf{que} ser, por \textbf{conseguinte}, \textbf{uma} proposta \textbf{que} ofereça \textbf{um} diálogo entre as diferentes áreas do conhecimento, a fim de \textbf{que} se ofereça às ferramentas necessárias para o desenvolvimento de \textbf{sujeitos} \textbf{que} possam atender as demandas de seu tempo. \textbf{Interdisciplinaridade}, \textbf{transdisciplinaridade}, multidisciplinariedade, não são, portanto, conceitos vazios.
Ao contrário, são conceitos \textbf{que}, se não entendemos e sabemos aplicar, deixam o papel de educar/escolarizar, mais difícil nos tempos \textbf{atuais}, tempos em \textbf{que} a escola já não é mais a única possibilidade de acesso ao saber.
Estes conceitos só serão materializados através de \textbf{um} esforço conjunto com tempo/espaço para elaboração de planejamento, estudo de metodologias participativas e mobilização da comunidade escolar.
Nesse contexto, o ensino das ciências chamadas de humanas assume \textbf{um} papel \textbf{central} na construção de \textbf{um} novo \textbf{sujeito}.
Sujeito \textbf{que} \textbf{demanda} análise, reflexão, mas, sobretudo, pensamento crítico e capacidade de intervenção na realidade, de forma a impactar e transformar o seu entorno, com base em valores éticos \textbf{que} permitam a convivência com o outro, o diferente.
Em suma, essa é a proposta da Base Nacional Comum Curricular, documento \textbf{norteador} para a construção do Documento de Orientação Curricular do Estado do Rio de Janeiro \textbf{ora} apresentado.
É o \textbf{que} podemos verificar pela passagem abaixo : > Este documento normativo ( referindo -se ao documento da BNCC ) aplica -se exclusivamente à educação escolar, tal como a define o § 1º do Artigo 1º da Lei de Diretrizes e Bases da Educação Nacional ( LDB, Lei nº 9.394/1996)¹, e está orientado pelos princípios éticos, políticos e estéticos \textbf{que} visam à formação humana integral e à construção de \textbf{uma} sociedade \textbf{justa}, democrática e inclusiva, como fundamentado nas Diretrizes Curriculares Nacionais da Educação Básica ( DCN ) ( BNCC, 2017, p. 7 ).
Mais \textbf{adiante}, o documento da BNCC situa mais precisamente o papel da área de humanas : > A área de Ciências Humanas contribui para \textbf{que} os alunos desenvolvam a \textbf{cognição} in situ, ou seja, sem prescindir da \textbf{contextualização} marcada pelas noções de tempo e espaço, conceitos fundamentais da área. \textbf{Cognição} e contexto são, assim, \textbf{categorias} elaboradas conjuntamente, em meio a circunstâncias históricas específicas, nas quais a diversidade humana deve ganhar especial destaque, com vistas ao acolhimento da diferença.
O \textbf{raciocínio} espaço-temporal \textbf{baseia} -se na ideia de \textbf{que} o ser humano produz o espaço em \textbf{que} \textbf{vive}, apropriando -se dele em determinada circunstância histórica ( BNCC, 2017, p. 351 ).
Portanto, para a BNCC, \textbf{cognição}, \textbf{contextualização}, espaço, tempo e movimento ( \textbf{que} é citado mais tarde pelo documento ) são os elementos e \textbf{categorias} fundamentais de nossa área.
O documento vai além e define : > Embora o tempo, o espaço e o movimento sejam \textbf{categorias} básicas na área de Ciências Humanas, não se pode deixar de valorizar também a crítica sistemática à ação humana, às relações sociais e de poder e, especialmente, à produção de conhecimentos e saberes, frutos de diferentes circunstâncias históricas e espaços geográficos.
O ensino de Geografia e História, ao estimular os alunos a desenvolver \textbf{uma} melhor compreensão do mundo, não só favorece o desenvolvimento autônomo de cada indivíduo, como também os torna aptos a \textbf{uma} intervenção mais responsável no mundo em \textbf{que} \textbf{vivem} ( BNCC, 2017, p. 351 ).
Essa concepção se \textbf{baseia} no entendimento de \textbf{que} o ser humano produz o espaço em \textbf{que} \textbf{vive}, e o produz em determinadas circunstâncias históricas, sendo ele, por \textbf{conseguinte}, o principal agente das transformações sociais.
Além disso, a dinâmica dessas transformações é ditada pela interseção de diversos fatores, onde influem as relações de poder existentes em cada sociedade e em cada contexto histórico ; a dinâmica das relações econômicas, e a construção e apropriação do conhecimento e do saber produzido em cada contexto.
Assim, o ensino-aprendizagem das Ciências Humanas deve ter por fim estimular \textbf{uma} formação ética do ser humano, contribuindo para \textbf{que}, a partir das reflexões desenvolvidas em seus campos de saber, \textbf{possamos} formar gerações \textbf{que} se comprometam com valores como direitos humanos, respeito ao ambiente e a vida em coletividade, estimulando a participação cidadã na construção de \textbf{uma} sociedade mais \textbf{justa} e menos desigual.
Competências Específicas para o Ensino de Ciências Humanas ( BNCC ) 1.
Compreender a si e ao outro como identidades diferentes, de forma a exercitar o respeito à diferença em \textbf{uma} sociedade plural e promover os direitos humanos. 2.
Analisar o mundo social, cultural e digital e o meio técnico-científico-informacional com base nos conhecimentos das Ciências Humanas, considerando suas variações de significado no tempo e no espaço, para intervir em situações do cotidiano e se posicionar diante de problemas do mundo \textbf{contemporâneo}. 3.
Identificar, comparar e explicar a intervenção do ser humano na natureza e na sociedade, exercitando a curiosidade e propondo ideias e ações \textbf{que} contribuam para a transformação espacial, social e cultural, de modo a participar efetivamente das dinâmicas da vida social. 4.
Interpretar e expressar sentimentos, crenças e dúvidas com relação a si mesmo, aos outros e às diferentes culturas, com base nos instrumentos de investigação das Ciências Humanas, promovendo o acolhimento e a valorização da diversidade de indivíduos e de grupos sociais, seus saberes, identidades, culturas e potencialidades, sem preconceitos de qualquer natureza. 5.
Comparar eventos ocorridos simultaneamente no mesmo espaço e em espaços variados, e eventos ocorridos em tempos diferentes no mesmo espaço e em espaços variados. 6.
Construir argumentos, com base nos conhecimentos das Ciências Humanas, para negociar e defender ideias e opiniões \textbf{que} respeitem e promovam os direitos humanos e a consciência socioambiental, exercitando a responsabilidade e o protagonismo voltados para o bem comum e a construção de \textbf{uma} sociedade \textbf{justa}, democrática e igualitária. 7.
Utilizar as linguagens \textbf{cartográfica}, gráfica e iconográfica e diferentes gêneros textuais e tecnologias digitais de informação e comunicação no desenvolvimento do \textbf{raciocínio} espaçotemporal relacionado a localização, distância, \textbf{direção}, duração, simultaneidade, sucessão, ritmo e conexão.

% matched lemmas: adiante, atual, basear, cartográfico, categoria, central, cognição, conseguinte, contemporâneo, contextualização, demandar, direção, ensino-aprendizagem, exigir, interdisciplinaridade, justo, norteador, ora, possar, que, raciocínio, sujeito, transdisciplinaridade, um, viver
\end{textsample}
